%%%%%%%%%%%%%%%%%%%%%%%%%%%%%%%%%%
% Página de resumen del proyecto %
%%%%%%%%%%%%%%%%%%%%%%%%%%%%%%%%%%
% !TEX root = A0.MiTFG.tex

\pagestyle{fancy}
\renewcommand{\headrulewidth}{0pt}
\addstarredchapter{Resumen}

\begin{center}
	\scshape
	E.T.S. de Ingeniería de Telecomunicación, Universidad de Málaga
\end{center}

\bigskip

\begin{center}
	\Large \scshape
	\textbf{\tfgtitlename}
\end{center}

\bigskip \bigskip \bigskip

\begin{minipage}{\textwidth}

\textbf{Autor:} \tfgauthorname


\medskip

\textbf{Tutor:} \tfgtutorname



\medskip

\textbf{Departamento:} Ingeniería de Comunicaciones

\medskip

\textbf{Titulación:} Grado en Ingeniería Telemática

\medskip

\textbf{Palabras clave:} comunicaciones ópticas, pulsos ultracortos, dispersión cromática, turbulencia atmosférica, turbulencia oceánica, propagación en medios no guiados, Matlab.

\bigskip \bigskip


\end{minipage}

\begin{center}
	\textbf{Resumen}
\end{center}

En este proyecto se lleva a cabo un estudio teórico y de simulación sobre los efectos de ensanchado y distorsión en la propagación de pulsos ópticos ultracortos debido a la dispersión cromática y a la turbulencia en medios atmosféricos y oceánicos. Las comunicaciones ópticas han despertado gran interés al ofrecer tasas de transmisión muy superiores a las tecnologías acústicas o de radiofrecuencia, aunque su aplicación en medios no guiados aún presenta retos, especialmente en el entorno oceánico.
Mediante simulaciones en Matlab se analiza el impacto de la dispersión cromática, que provoca el ensanchamiento temporal de los pulsos y limita la capacidad de transmisión y de la turbulencia, que introduce fluctuaciones en fase y amplitud que reducen la potencia recogida en el receptor. Los resultados obtenidos permiten comprender mejor estas limitaciones y su influencia en el diseño de sistemas de comunicación óptica en entornos atmosféricos y submarinos. 



\blankpage
